% Cote 119: offprints, reports, courses, CV and letters (typed and printed part): archive cote 119, site page 296 (PDF page 297).
% typescript. Draft compared: none.
% AI-transcribed by Claude Opus 5.5 (claude-opus-5-5, adaptive thinking, effort high), Message Batches API, batch r119_full (msgbatch_01DwZNm5Y7YsKBWyAN4sFs8w), 2026-11-13; not yet checked by eye.
\scanpage{119/296}
\mspage{- 9 -}

Le cas où le rang $n$ de $E$ est égal à $6$ pose également des problèmes
intéressants. Notons $e : \mathrm{Gal}(k_s/k) \to S_6$ l'homomorphisme qui correspond à
$E$ par la théorie de Galois. En composant $e$ avec un automorphisme extérieur
de $S_6$ on obtient un homomorphisme $e' : \mathrm{Gal}(k_s/k) \to S_6$ qui correspond
à une autre algèbre étale $E'$ de rang $6$ ( `` résolvante sextique '' ).
\uline{Peut-on déterminer} $q_{E'}$ \uline{à partir de} $q_E$ ? C'est vrai lorsque $w_{\hand{1}}(q_E) = 0$ ,
autrement dit lorsque les images de $e$ et $e'$ sont contenues dans le
groupe alterné $A_6$ ; on peut en effet prouver que l'on a dans ce cas
$q_{E'} \hand{$\simeq$} 2q_E$ (mais pas $q_{E'} \hand{$\simeq$} q_E$ en général, bien que $q_E$ et $q_{E'}$ aient les
mêmes invariants $w_1$ et $w_2$ ). Lorsque l'on a à la fois $w_1(q_E) = 0$
et $w_2(q_E) = 0$, on peut se demander si $q_E$ est isomorphe à la forme
unité $\hand{$\langle$}1,1,\ldots,1\hand{$\rangle$}$. C'est vrai si $k$ est un corps de nombres (ou un corps
de fonctions rationnelles sur un corps de nombres) ; c'est probablement
faux en général, mais je n'ai pas réussi à construire un contre-exemple.

\textbf{§6. La théorie de Bayer-Lenstra : les bases normales autoduales}

Soit $G$ un groupe fini. On s'intéresse aux G-\uline{algèbres galoisiennes}
sur $k$, ou, ce qui revient au même, aux G-\uline{torseurs} sur $k$, $G$ étant considéré
comme un groupe algébrique de dimension $0$ sur $k$.Une telle algèbre $L$ est
déterminée, à isomorphisme (non unique) près, par la donnée d'un homomorphisme
continu
\[ \hand{$\varphi$}_L : \mathrm{Gal}(k_s/k) \to G. \]
Lorsque $\hand{$\varphi$}_L$ est surjectif, $L$ est un corps, et c'est une extension
galoisienne de $k$ de groupe de Galois isomorphe à $G$.

Dans un travail récent (\uline{Amer}.\uline{J}.\uline{Math}. 112, 1990), E.Bayer et H.Lenstra
s'intéressent au cas où $L$ possède une \uline{base normale autoduale} (`` BNA '');
cela signifie qu'il existe un élément $x$ de $L$ tel que $q_L(x) = 1$ et
que $x$ soit orthogonal (relativement à $q_L$)à tous les $gx$, $g \hand{$\in$} G$, $g \hand{$\neq$} 1$ .
(Ainsi, les $gx$ forment une ``base normale '' de $L$, et cette base est sa
propre duale relativement à $q_L$ .)

On peut donner un critère cohomologique pour l'existence d'une BNA : si
$U_G$ désigne le groupe unitaire de l'algèbre à involution $k\hand{$[$}G\hand{$]$}$ , on a un
plongement canonique de $G$ dans $U_G(k)$; en composant $\hand{$\varphi$}_L$ avec ce plongement
on obtient un homomorphisme $\mathrm{Gal}(k_s/k) \to U_G(k)$ , homomorphisme que l'on peut
regarder comme un 1-cocycle de $\mathrm{Gal}(k_s/k)$ à valeurs dans $U_G(k_s)$. La classe
$\hand{$\varepsilon$}_L$ de ce cocycle est un élément de $H^1(k, U_G)$ . \uline{On a} $\hand{$\varepsilon$}_L = 0$ \uline{si et}
\uline{seulement si} $L$ \uline{a une} BNA .

\margin{296}
